% GENERATED FILE. Edit canonical lessons, then run npm run generate:books.
% canonical-source-hash: fnv1a64:bd705c4fa3783f58
% canonical-lessons: GE-C379-bewerten, GE-C379-bewertung, GE-C379-benoten, GE-C379-nachholen, GE-C379-korrigieren

\chapter{Noten und Bewertungen}
\label{ch:ge-noten-und-bewertungen}
\hlchaptermodality{379}

\begin{chapteropening}
\textbf{By the end of this chapter:} \emph{I can talk about how work is rated and graded, correcting mistakes, and making up a missed exam.}

Complete the last lesson of chapter 379: I can talk about how work is rated and graded, correcting mistakes, and making up a missed exam.
\end{chapteropening}

\section[bewerten]{bewerten — to rate, to evaluate}

\label{lesson:GE-C379-bewerten}



\begin{quote}
Before the new one: say the German for the summary, then the German for to sum up.
\end{quote}

\subsection*{You'll want to know: bewerten}
\textbf{bewerten} — \textquotedblleft{}to rate, to evaluate\textquotedblright{}.

\begin{quote}
Die Lehrer bewerten die Aufsätze nach klaren Kriterien.
\end{quote}

\begin{quote}
\emph{The teachers evaluate the essays by clear criteria.}
\end{quote}

Inseparable: \emph{sie bewertet, sie hat bewertet}. Online: \emph{ein Hotel bewerten} — to rate a hotel.

\begin{cousinweb}
\emph{bewerten} is built with \emph{be-} on \emph{werten} \textquotedblleft{}to value\textquotedblright{}, from \emph{Wert} \textquotedblleft{}worth\textquotedblright{}, inherited from Old High German \emph{werd} and Proto-Germanic \emph{werþaz} (EtymDB), a cousin of English \emph{worth}.
\end{cousinweb}

\begin{grammarlens}[title={What you've built}]
One more word for telling, explaining and describing at B1.
\end{grammarlens}

\subsection*{Your turn}
\begin{itemize}
\raggedright
  \item \emph{Say it:} \emph{bewerten}
  \item \emph{Say it:} \emph{bewerten}, once more
  \item \emph{Say it:} a sentence of your own with \emph{bewerten}
  \item \emph{Recall:} say the German for the summary, then the German for to sum up, then say \emph{bewerten} again
\end{itemize}

\begin{tcolorbox}[breakable,colback=teal!4,colframe=teal!35!black,title={Before you move on}]
What does bewerten mean? (\textquotedblleft{}To rate, to evaluate\textquotedblright{}.) Say it once more.
\end{tcolorbox}
\section[die Bewertung]{die Bewertung — the evaluation, the rating; the review}

\label{lesson:GE-C379-bewertung}



\begin{quote}
Before the new one: say the German for to sum up, then the German for to rate.
\end{quote}

\subsection*{You'll want to know: die Bewertung}
\textbf{die Bewertung} — \textquotedblleft{}the evaluation, the rating; the review\textquotedblright{}.

\begin{quote}
Das Restaurant hat im Internet sehr gute Bewertungen.
\end{quote}

\begin{quote}
\emph{The restaurant has very good reviews online.}
\end{quote}

Plural: \emph{die Bewertungen}. At school, the marks and the comments together.

\begin{cousinweb}
\emph{Bewertung} is the noun of \emph{bewerten}, built on \emph{Wert} \textquotedblleft{}worth\textquotedblright{}, Proto-Germanic \emph{werþaz} (EtymDB), with \emph{-ung}.
\end{cousinweb}

\begin{grammarlens}[title={What you've built}]
One more word for telling, explaining and describing at B1.
\end{grammarlens}

\subsection*{Your turn}
\begin{itemize}
\raggedright
  \item \emph{Say it:} \emph{die Bewertung}
  \item \emph{Say it:} \emph{die Bewertung}, once more
  \item \emph{Say it:} a sentence of your own with \emph{die Bewertung}
  \item \emph{Recall:} say the German for to sum up, then the German for to rate, then say \emph{die Bewertung} again
\end{itemize}

\begin{tcolorbox}[breakable,colback=teal!4,colframe=teal!35!black,title={Before you move on}]
What does die Bewertung mean? (\textquotedblleft{}The evaluation, the rating; the review\textquotedblright{}.) Say it once more.
\end{tcolorbox}
\section[benoten]{benoten — to grade (with a mark)}

\label{lesson:GE-C379-benoten}



\begin{quote}
Before the new one: say the German for to rate, then the German for the evaluation.
\end{quote}

\subsection*{You'll want to know: benoten}
\textbf{benoten} — \textquotedblleft{}to grade (with a mark)\textquotedblright{}.

\begin{quote}
Die Klausuren werden erst nächste Woche benotet.
\end{quote}

\begin{quote}
\emph{The exams will not be graded until next week.}
\end{quote}

Inseparable: \emph{er benotet, er hat benotet}. From \emph{die Note} — the mark.

\begin{cousinweb}
\emph{benoten} is built with \emph{be-} on \emph{Note} \textquotedblleft{}mark, grade\textquotedblright{}, from Latin \emph{nota} \textquotedblleft{}a mark, a sign, a note\textquotedblright{} (EtymDB; Lewis \& Short).
\end{cousinweb}

\begin{grammarlens}[title={What you've built}]
One more word for telling, explaining and describing at B1.
\end{grammarlens}

\subsection*{Your turn}
\begin{itemize}
\raggedright
  \item \emph{Say it:} \emph{benoten}
  \item \emph{Say it:} \emph{benoten}, once more
  \item \emph{Say it:} a sentence of your own with \emph{benoten}
  \item \emph{Recall:} say the German for to rate, then the German for the evaluation, then say \emph{benoten} again
\end{itemize}

\begin{tcolorbox}[breakable,colback=teal!4,colframe=teal!35!black,title={Before you move on}]
What does benoten mean? (\textquotedblleft{}To grade (with a mark)\textquotedblright{}.) Say it once more.
\end{tcolorbox}
\section[nachholen]{nachholen — to make up (for), to catch up on}

\label{lesson:GE-C379-nachholen}



\begin{quote}
Before the new one: say the German for the evaluation, then the German for to grade.
\end{quote}

\subsection*{You'll want to know: nachholen}
\textbf{nachholen} — \textquotedblleft{}to make up (for), to catch up on\textquotedblright{}.

\begin{quote}
Er war krank und muss die Prüfung nächste Woche nachholen.
\end{quote}

\begin{quote}
\emph{He was ill and has to make up the exam next week.}
\end{quote}

Separable: \emph{er holt nach}, \emph{er hat nachgeholt}.

\begin{cousinweb}
\emph{nachholen} is built from \emph{nach} \textquotedblleft{}after\textquotedblright{} and \emph{holen} \textquotedblleft{}to fetch\textquotedblright{}, from Middle High German \emph{holen} (EtymDB): you fetch what you missed afterwards.
\end{cousinweb}

\begin{grammarlens}[title={What you've built}]
One more word for telling, explaining and describing at B1.
\end{grammarlens}

\subsection*{Your turn}
\begin{itemize}
\raggedright
  \item \emph{Say it:} \emph{nachholen}
  \item \emph{Say it:} \emph{nachholen}, once more
  \item \emph{Say it:} a sentence of your own with \emph{nachholen}
  \item \emph{Recall:} say the German for the evaluation, then the German for to grade, then say \emph{nachholen} again
\end{itemize}

\begin{tcolorbox}[breakable,colback=teal!4,colframe=teal!35!black,title={Before you move on}]
What does nachholen mean? (\textquotedblleft{}To make up (for), to catch up on\textquotedblright{}.) Say it once more.
\end{tcolorbox}
\section[korrigieren]{korrigieren — to correct}

\label{lesson:GE-C379-korrigieren}



\begin{quote}
Before the new one: say the German for to grade, then the German for to make up.
\end{quote}

\subsection*{You'll want to know: korrigieren}
\textbf{korrigieren} — \textquotedblleft{}to correct\textquotedblright{}.

\begin{quote}
Die Lehrerin korrigiert jeden Abend Hefte.
\end{quote}

\begin{quote}
\emph{The teacher corrects exercise books every evening.}
\end{quote}

Perfekt: \emph{hat korrigiert}. \emph{Sich korrigieren} — to correct yourself.

\begin{cousinweb}
\emph{korrigieren} is a learned borrowing from Latin \emph{corrigere} \textquotedblleft{}to make straight, to set right\textquotedblright{} (Lewis \& Short), built from \emph{con-} and \emph{regere} \textquotedblleft{}to guide, to rule\textquotedblright{}.
\end{cousinweb}

\begin{grammarlens}[title={What you've built}]
One more word for telling, explaining and describing at B1.
\end{grammarlens}

\subsection*{Your turn}
\begin{itemize}
\raggedright
  \item \emph{Say it:} \emph{korrigieren}
  \item \emph{Say it:} \emph{korrigieren}, once more
  \item \emph{Say it:} a sentence of your own with \emph{korrigieren}
  \item \emph{Recall:} say the German for to grade, then the German for to make up, then say \emph{korrigieren} again
\end{itemize}

\begin{tcolorbox}[breakable,colback=teal!4,colframe=teal!35!black,title={Before you move on}]
What does korrigieren mean? (\textquotedblleft{}To correct\textquotedblright{}.) Say it once more.
\end{tcolorbox}
